\documentclass[a4paper,11pt]{article}
\usepackage[utf8]{inputenc}
\usepackage[T1]{fontenc}
\usepackage[english]{babel}
\usepackage{geometry}
\geometry{left=2cm,right=2cm,top=2cm,bottom=2cm}

\usepackage{lmodern}
\usepackage{xcolor}
\definecolor{titleblue}{RGB}{0,51,140}
\usepackage{hyperref}
\usepackage{titlesec}

\titleformat{\section}{\color{titleblue}\large\bfseries}{\thesection}{1em}{}

\begin{document}

\begin{center}
    {\LARGE\bfseries\color{titleblue} Research Statement}\\[6pt]
    {\large Distributed Abductive Reasoning for Self-Explaining Digital Twins}\\[4pt]
    \textbf{Ismail AISSAOUI} $|$ State Engineer in AI
\end{center}

\bigskip

\section{Introduction and Motivation}
As digital twins transition from centralized simulations to distributed cyber-physical companions, they must handle the "reality gap" between model predictions and real-world observations. Current systems often treat these discrepancies as numerical errors to be minimized, without providing a human-interpretable explanation of their root causes. This research statement proposes a hierarchical and distributed explanation architecture based on abductive reasoning, enabling digital twins to interpret discrepancies as actionable knowledge about the system state and its requirements.

\section{Current Research: Physics-Informed Digital Twins}
In my current work at Sonatrach, I have developed a generative surrogate for gas turbine monitoring using \textbf{Mamba-SSM} and \textbf{xLSTM} architectures. A key achievement of this work is the deployment of these complex hybrid models onto \textbf{Edge devices}, where real-time performance and local autonomy are critical. However, my experience has shown that local model adaptation is insufficient if the digital twin cannot "explain" why a mismatch occurred—whether due to sensor drift, environmental shifts, or physical degradation. My background in edge-deployed AI provides a solid foundation for building the distributed reasoning layers needed for explainable twins.

\section{Proposed Research Directions}
Building upon the objectives of the \textbf{Distributed AI} project at Université Côte d’Azur, I propose to investigate the following axes:

\subsection{1. Abductive Reasoning for Hybrid Models}
Unlike deduction or induction, abductive reasoning focuses on generating the most plausible explanations for unexpected observations. I propose to develop abductive mechanisms that can bridge the gap between physics-based equations and data-driven learning. By formalizing the "surprise" encountered by a digital twin, we can generate hypotheses that distinguish between internal model deficiencies and external physical disturbances.

\subsection{2. Hierarchical Edge-Fog-Cloud Explanation Architecture}
I intend to investigate a distributed reasoning framework where low-level anomalies are detected at the \textbf{Edge}, refined into local explanations at the \textbf{Fog} level, and finally validated against global system requirements in the \textbf{Cloud}. This multi-level approach ensures that explanations are contextual and traceable, providing stakeholders with different granularities of insight depending on their roles.

\subsection{3. Requirement-Aware Self-Explanation}
A critical challenge in self-explaining systems is ensuring that explanations remain consistent with the broader system requirements. I propose to integrate "requirement-aware" constraints into the abductive process, ensuring that the hypotheses generated by the digital twin do not violate safety or performance boundaries.

\section{Alignment with the EDT Program}
My background in distributed AI and industrial implementation aligns with the goal of creating a national platform for digital twin engineering. I am eager to contribute to the **Artemis platform** by providing a distributed reasoning layer that enhances the transparency and trustworthiness of hybrid digital twins across the consortium.

\section{Conclusion}
The goal of my PhD will be to move beyond reactive error correction towards proactive, self-explaining digital twins. By combining abductive reasoning with distributed architectures, I aim to provide the foundational tools necessary for the next generation of transparent and trustworthy cyber-physical companions.

\end{document}
